\clearpage
\phantomsection%
\addcontentsline{toc}{chapter}{ABSTRACT}
\thispagestyle{fancy}

\begin{center}
	\large \bfseries \MakeUppercase{Abstract}\\
	\normalsize \normalfont {\thetitleEN}\\
	\normalsize \normalfont {\theauthor}\\
	\bigskip

	\normalsize \normalfont \justifying \singlespacing
	Indonesian Sign Language (SIBI) is the primary communication system for the Deaf community in Indonesia; however, the lack of automated video-based recognition systems remains a significant barrier to communication accessibility. This study analyzes the performance of the Video Masked Autoencoder (VideoMAE) model in classifying ten basic SIBI gesture classes using a dataset of 924 video clips collected directly at the State Special School (SLBN) PKK of Lampung Province. Evaluation was conducted using a Stratified 5-Fold Cross Validation scheme. Twelve VideoMAE experimental configurations (E1--E12) were systematically designed to investigate the effects of backbone architecture (ViT-Small, ViT-Base), pre-training dataset (Kinetics-400, SSV2), training strategy (two-stage fine-tune vs. scratch), frame sampling strategy (uniform vs. dense), number of epochs, and augmentation techniques (mixup, CutMix, reprob). ViViT-B/16x2 and I3D were also evaluated as comparative models. Results demonstrate that pre-trained weights are an absolute prerequisite: models trained from random initialization achieved only 11--12\% accuracy, whereas two-stage pre-training with SSV2 weights improved accuracy to 63.85\%. In the controlled sampling comparison, uniform sampling outperformed dense sampling (63.85\% vs.\ 56.07\%). The most significant finding, however, is that mixup, CutMix, and reprob augmentations proved counterproductive on small-scale sign language datasets because they disrupt the spatio-temporal context of gestures; disabling them raised accuracy to 93.94\%. The best overall configuration was VideoMAE ViT-Base with Kinetics-400 weights and dense sampling without augmentation (E12), achieving a mean accuracy of 97.18\%, indicating that sampling effectiveness depends on its interaction with the broader model configuration, and marginally surpassing I3D (96.86\%) and ViViT (94.81\%).\par

    \vspace{20pt}
	\raggedright\textbf{Keywords: VideoMAE, SIBI Sign Language, Video Classification, Vision Transformer, Fine-tuning}

	\vfill

\end{center}
\clearpage
